% Discussion snippet: FSDS shares as a prior for the next online-learning round.
% Intended to sit after Tables~\ref{tab:msrvtt-mm-shares}--\ref{tab:msrvtt-mm-pervideo-auc}.
% Requires amsmath (optional) and the labels already defined in
% MSRVTT_Multimodal_Attribution_tables_only.tex.

\paragraph{From FSDS attribution to the next online-learning round.}
The modality-specific FSDS analysis is not only a retrospective decomposition of covariate shift $P(X\mid W)$: it supplies an operational prior for the subsequent online-learning (OL) round.
Across RF-Domain, coordinate-MMD, and PO-risk VIMP, the within-video early-versus-late window shift is carried primarily by the visual block ($\pi_{\mathrm{v}}\approx 0.77$--$0.90$), with a nontrivial audio contribution ($\pi_{\mathrm{a}}\approx 0.10$--$0.19$) and a near-null text share (Table~\ref{tab:msrvtt-mm-shares}).
A Friedman test on per-video RF shares rejects equal modality contributions ($p=2.87\times 10^{-7}$), and Holm-adjusted Wilcoxon tests reject all three pairwise equalities at $p<10^{-4}$ (Table~\ref{tab:msrvtt-mm-pervideo-tests}).
Region-wise Cohen's $d$ reproduces the same ranking (mean $|d|$ of $0.787$/$0.518$/$0.000$ for video/audio/text).
Signed pooled $d$ cancels across videos ($0/24$ Holm hits at the pooled-region level), so the shift is \emph{video-heterogeneous} rather than a single global direction---which is the quantity a pooled OL gradient would otherwise average away.

These facts constrain the next OL protocol in three concrete ways.
\emph{(i)~Where to adapt.}
Update \emph{capacity and refresh} should be concentrated on the visual encoder and, secondarily, the audio encoder.
That is not the same as a larger SGD step: a large covariate intensity on those towers \emph{lowers} $\eta_m$.
Clip-level caption embeddings are window-invariant on this extract, hence a \emph{structural} zero for $W$-shift rather than evidence that language is uninformative for retrieval; a time-aligned ASR stream would be required before attributing residual $W$-shift to text.
\emph{(ii)~How to monitor.}
Drift triggers should be video-clustered (the same sampling unit as the bootstrap and Friedman tests).
A global MMD or pooled $d$ will miss the heterogeneous $|d|$ that FSDS detects.
\emph{(iii)~How to fuse.}
The RF shares provide a hot-start for modality gates, with a documented per-video exception (video~13 is audio-dominant; Table~\ref{tab:msrvtt-mm-pervideo-auc}) that argues against a frozen global fusion weight.
Until a window-aligned text encoder is in the stream, a rising text VIMP in OL is a diagnostic for leakage, not a reason to retune the language tower against $W$.

The per-head learning-rate map is \emph{not} $\eta_m=\eta_0\pi_m^{\mathrm{RF}}$.
FSDS $\pi$ and $c_m^{\mathrm{cov}}$ are covariate shift $P(X^{(m)}\mid W)$.
A large covariate intensity on modality $m$ \emph{lowers} $\eta_m$: the embedding moved, so a large step would chase $X$ and fit $W$ as a shortcut.
A large concept drift $\delta_m^{\mathrm{concept}}$ in $P(Y\mid X^{(m)})$ \emph{raises} $\eta_m$.
Both quiet (clip-level text on this $W$) $\Rightarrow$ $\eta_{\mathrm{t}}\approx 0$; do not invert $\pi$ or text inherits the largest step.
Refresh / KV invalidation still follows covariate mass; that is a memory schedule, not the SGD sign.
Ranking of $\pi$ can be shared across batches; the scale of $c_m$ and $\delta_m$ can be session-specific (video~13).
Cosine lag and round-to-round VIMP comparison are separate boards.

In short, the offline FSDS readout is the design matrix for the next round of streaming results: re-encode vision and audio, cluster by video, treat clip-level text as a negative control, shrink $\eta_m$ when covariate shift is large, and raise $\eta_m$ only when concept drift is large.
